\hypertarget{enums-and-topological-sorts-enum-graphlib}{%
\chapter{\texorpdfstring{Enums and Topological Sorts (\texttt{enum},
\texttt{graphlib})}{Enums and Topological Sorts (enum, graphlib)}}\label{enums-and-topological-sorts-enum-graphlib}}

This chapter explores advanced data categorization and dependency
resolution modules that leverage Python's metaclassing and algorithmic
strengths.

\hypertarget{enum-metaprogramming-constant-mappings}{%
\subsection{\texorpdfstring{47.1 \texttt{enum}: Metaprogramming Constant
Mappings}{47.1 enum: Metaprogramming Constant Mappings}}\label{enum-metaprogramming-constant-mappings}}

As introduced in Chapter 10, \passthrough{\lstinline!enum!} is more than
a simple constant list.

\hypertarget{intenum-and-intflag}{%
\subsubsection{\texorpdfstring{1. \texttt{IntEnum} and
\texttt{IntFlag}}{1. IntEnum and IntFlag}}\label{intenum-and-intflag}}

\begin{itemize}
\tightlist
\item
  \textbf{\passthrough{\lstinline!IntEnum!}}: Subclasses both
  \passthrough{\lstinline!int!} and \passthrough{\lstinline!Enum!}. It
  allows comparisons with raw integers
  (\passthrough{\lstinline!Color.RED == 1!}).
\item
  \textbf{\passthrough{\lstinline!IntFlag!}}: Supports bitwise
  operations (\passthrough{\lstinline!|!}, \passthrough{\lstinline!\&!},
  \passthrough{\lstinline!\^!}, \passthrough{\lstinline!\~!}). It is
  useful for representing hardware registers or permission bitmasks.
\item
  \textbf{Internals}: \passthrough{\lstinline!IntFlag!} members are
  combined using a specialized version of the
  \passthrough{\lstinline!EnumMeta!} metaclass that ensures bitwise
  results are still valid members of the Flag class.
\end{itemize}

\hypertarget{the-auto-helper}{%
\subsubsection{\texorpdfstring{2. The \texttt{auto()}
Helper}{2. The auto() Helper}}\label{the-auto-helper}}

\passthrough{\lstinline!auto()!} is a sentinel object. During class
construction, the metaclass detects \passthrough{\lstinline!auto()!} and
assigns an appropriate value (usually an incrementing integer).

\hypertarget{graphlib-dependency-resolution}{%
\subsection{\texorpdfstring{47.2 \texttt{graphlib}: Dependency
Resolution}{47.2 graphlib: Dependency Resolution}}\label{graphlib-dependency-resolution}}

Python 3.9 introduced \passthrough{\lstinline!graphlib!} to provide a
standard way to perform topological sorting of graphs.

\hypertarget{topological-sorting}{%
\subsubsection{1. Topological Sorting}\label{topological-sorting}}

A topological sort is a linear ordering of vertices such that for every
directed edge \((u, v)\), \(u\) comes before \(v\). This is the
foundation of build systems (e.g., \passthrough{\lstinline!make!},
\passthrough{\lstinline!ninja!}) and task schedulers.

\hypertarget{topologicalsorter-internals}{%
\subsubsection{\texorpdfstring{2. \texttt{TopologicalSorter}
Internals}{2. TopologicalSorter Internals}}\label{topologicalsorter-internals}}

The \passthrough{\lstinline!TopologicalSorter!} class implements a
\textbf{Kahn's Algorithm} variant. 1. It calculates the ``in-degree''
(number of incoming edges) for every node. 2. It maintains a queue of
nodes with in-degree 0 (those with no dependencies). 3. As nodes are
``prepared'' and ``done'', it decrements the in-degree of their
neighbors, adding new 0-degree nodes to the queue. 4. \textbf{Cycle
Detection}: If the graph is exhausted but nodes remain with in-degree
\textgreater{} 0, it raises a \passthrough{\lstinline!CycleError!}.

\hypertarget{parallel-execution}{%
\subsubsection{3. Parallel Execution}\label{parallel-execution}}

\passthrough{\lstinline!graphlib!} is designed for parallel runners. You
can call \passthrough{\lstinline!get\_ready()!} to get all nodes that
can be executed immediately, work on them in separate threads/processes,
and call \passthrough{\lstinline!done()!} as they finish to unlock the
next tier of dependencies.

\begin{center}\rule{0.5\linewidth}{0.5pt}\end{center}
